%% kapa: exam=1 round=33 year=2022 period=1 form=A subject=economics subject_ko=경제학원론 seq=36 number=76
%% source: Q-Net 감정평가사 기출문제 (cst00309) artlSeq=5209997 file=raw/1008/5209997/2022년도 제33회 감정평가사 1차시험 1교시(민법,경제학,부동산학).pdf
%% conversion: kordoc 4.24.51 -> tools/gen1.py (auto); review=proofread-sonnet; flags=none
\begin{question}{76}
아래의 폐쇄경제 IS-LM 모형에서 도출된 총수요곡선으로 옳은 것은? (단, $r$은 이자율, $Y$는 국민소득, $M^d$는 명목화폐수요량, $P$는 물가수준, $M^s$는 명목화폐공급량이고, $Y>20$이다.)
\begin{framedtext}
○ IS곡선: $r = 10 - 0.4Y$ \par
○ 실질화폐수요함수: $\dfrac{M^d}{P} = 0.1Y - r$ \par
○ 명목화폐공급함수: $M^s = 4$
\end{framedtext}
\begin{choices}
  \item $P = \dfrac{1}{2(Y-20)}$
  \item $P = \dfrac{1}{Y-20}$
  \item $P = \dfrac{2}{Y-20}$
  \item $P = \dfrac{4}{Y-20}$
  \item $P = \dfrac{8}{Y-20}$
\end{choices}
\end{question}
\answer{5}
